\documentclass[UTF8,fontset=none,11pt,a4paper]{ctexart}
\usepackage[margin=2.5cm]{geometry}
\usepackage{amsmath,amssymb,bm,booktabs,tabularx}
\usepackage{xurl}
\usepackage[hidelinks]{hyperref}
\setCJKmainfont[AutoFakeBold=2]{AR PL UMing CN}
\setCJKsansfont{WenQuanYi Micro Hei}
\setmonofont{DejaVu Sans Mono}
\DeclareMathOperator{\Tr}{Tr}
\newcommand{\ii}{\mathrm{i}}
\newcommand{\OPE}{\mathcal O}
\newcommand{\TE}{\widetilde F^{E}}
\newcommand{\TM}{\widetilde F^{M}}
\title{HYP 链接下的胶子 PDF 研究笔记\\[0.2em]\large 直 Wilson line 算符约定与原始 OPE}
\author{}
\date{2026 年 11 月 24 日}
\begin{document}
\maketitle

本笔记记录 4D HYP smearing 10 次的直 Wilson line 胶子 OPE。
分析采用 unpolarized 的正位移 \(\OPE_U^E(+|z|)\) 和 helicity 的奇差
\(\OPE_H^E(+z)-\OPE_H^E(-z)\)。两个通道的欧氏组合均为时间分量和减去
两倍空间分量。这里的 \(z\) 表示沿 Wilson line 的有符号空间位移。
\(\nu=P_z z\) 为 Ioffe time，计算时用自然单位。

\section{Minkowski 与 Euclidean 字典}

Minkowski 顺序为 \((0,1,2,3)=(t,x,y,z)\)，度规
\(g=\mathrm{diag}(1,-1,-1,-1)\)。取
\(\epsilon_M^{0123}=+1\)，降指标后 \(\epsilon^M_{0123}=-1\)。
欧氏格点的代码顺序为 \((x,y,z,t)=(0,1,2,3)\)，取
\(\epsilon^E_{xyzt}=+1\)。采用同一 Wick 旋转
\(x_M^0=-\ii x_E^t\)，场强字典为
\begin{align}
  F^M_{0i}&=\ii F^E_{ti}, & F^M_{ij}&=F^E_{ij},\\
  \TM_{\mu\nu}&=\tfrac12\epsilon^M_{\mu\nu\rho\sigma}F_M^{\rho\sigma},
    & \TE_{\mu\nu}&=\tfrac12\epsilon^E_{\mu\nu\rho\sigma}F^E_{\rho\sigma},\\
  \TM_{0i}&=\TE_{ti}, & \TM_{ij}&=-\ii\TE_{ij}.
\end{align}
例如，欧氏代码中
\(\TE_{30}=-F^E_{12}\)，
\(\TE_{31}=+F^E_{02}\)，
\(\TE_{01}=+F^E_{23}\)。
Minkowski 的对偶张量中使用 \(F_M^{\rho\sigma}\)，因此升降指标的符号必须
与 Wick 相位一起追踪。这里的 Wick 相位属于算符层。核子 helicity 投影、
动量与自旋方向和最终比值的复数相位在 \(C_3/C_2\) 层另行定义。
相关字典也见梯度流笔记第 8.3--8.7 节 \cite{FlowNote}。

\section{等时直线缩并与通道组装}

令 \(y=x+z\hat{\bm n}\)，\(W_z(x)\) 为从 \(x\) 到 \(y\) 的有向链接乘积。
两端场强与 Wilson line 取自同一套已经 HYP 平滑的规范场。
对每一个时间片 \(t/a\)，程序保存
\begin{equation}
 M^E_{A;B}(t,z)=\sum_{\bm x}\Tr\!\left[
    A(y,t)W_z^\dagger(x)B(x,t)W_z(x)\right] .
 \label{eq:bilocal}
\end{equation}
\(U_{\mu\nu}=M^E_{F_{\mu\nu};F_{\mu\nu}}\)，
\(H_{\mu\nu}=M^E_{F_{\mu\nu};\widetilde F_{\mu\nu}}\)。
颜色迹用基本表示。若两端均为无迹场强，取
\(\Tr(T^aT^b)=\delta^{ab}/2\)，则基本表示迹为相应伴随线颜色缩并的
\(1/2\)。当前格点内核使用四叶 Clover
\[
 F^{E,\mathrm{HYP}}_{\mu\nu}=-\frac{\ii}{8}
 (Q_{\mu\nu}-Q_{\mu\nu}^{\dagger})
\]
并保留色迹。与无迹 Clover 的差是
\(\tfrac13\sum_{\bm x}\Tr A(y,t)\Tr B(x,t)\)，已有双场强文件无法单独给出这项。
梯度流笔记的生产方案从 thin links 做 Wilson flow，并在两个端点去迹。
二者可以采用同一张量组合，分别保留各自的算符归一化和重整化定义。

\subsection{Unpolarized}

沿 \(z\) 方向时，Khan 等的 Minkowski 组合
\(\sum_{i=x,y}M^M_{0i;i0}+2M^M_{xy;yx}\) 消去了 \(\mathcal M_{gg}\)
\cite{Khan,BMRU}。Wick 旋转后
\(F^M_{0i}=+\ii F^E_{ti}\)、\(F^M_{i0}=-\ii F^E_{ti}\)，两项相乘为正，
而空间反序给出负号。这给出 Khan 等的式 (19) 以及 Chen 等
的式 (2)--(3) \cite{Khan,Chen}
\begin{equation}
 \OPE_U^E(z)=U_{tx}(z)+U_{ty}(z)-2U_{xy}(z)=T_U(z)-S_U(z).
 \label{eq:unpol}
\end{equation}
其中 \(T_U=U_{tx}+U_{ty}\)，\(S_U=2U_{xy}\)。空间系数 2 来自
\((i,j)=(x,y),(y,x)\) 两个反序分量。

完整周期空间求和下，把 \(-z\) 路径的起点移到另一端并循环颜色迹，
得到
\[
 M^E_{A;B}(-z)=M^E_{B;A}(+z),\qquad
 U_{\mu\nu}(-z)=U_{\mu\nu}(+z).
\]
这个等式在每个组态、每个时间片成立。实际读入时取 \(+|z|\)，也包括局域
\(z=0\)。旧偶和等于 \(2\OPE_U^E(+|z|)\)，是历史存储的归一化。
若用正半轴实现全轴 Fourier 变换，\(1/(2\pi)\) 的全轴变换系数对应
\(1/\pi\) 的正半轴余弦变换系数。

\subsection{Helicity}

Egerer 等式 (5) 定义奇差，式 (10) 的 Minkowski 组合为
\(\sum_{i=x,y}\widetilde M^M_{0i;0i}+2\widetilde M^M_{xy;xy}\)。
Balitsky、Morris、Radyushkin 的式 (2.17)--(2.21) 说明它消去
\(\widetilde{\mathcal M}_{gg}\)，并保留有限动量的
\(m_N^2/P_z^2\) 项 \cite{Egerer,BMRH}。时间项由
\(F^M_{0i}=\ii F^E_{ti}\) 带来 \(+\ii\)。空间项中的普通 \(F_{xy}\)
在 Wick 旋转下保持不变，对偶场强给出
\begin{equation}
 \TM_{xy}
 =\epsilon^M_{xy0z}F_M^{0z}
 =(-1)(-\ii F^E_{tz})
 =-\ii\TE_{xy},\qquad \TE_{xy}=F^E_{zt}.
 \label{eq:helicity-wick-xy}
\end{equation}
因此 Minkowski 的加号组合是
\(\ii(H_{tx}+H_{ty}-2H_{xy})\)。按代码只保存欧氏缩并的约定，定义
\begin{align}
 \OPE_H^E(z)&=H_{tx}(z)+H_{ty}(z)-2H_{xy}(z)=T_H(z)-S_H(z),\\
 \OPE_{H,\mathrm{odd}}^E(z)&=\OPE_H^E(+z)-\OPE_H^E(-z).
 \label{eq:hel}
\end{align}
这里 \(T_H=H_{tx}+H_{ty}\)，\(S_H=2H_{xy}\)。用
\(K_{\mu\nu;\rho\sigma}=M^E_{F_{\mu\nu};F_{\rho\sigma}}\) 展开，两个空间分量为
\begin{align}
 U_{xy}&=K_{xy;xy},    & -2U_{xy}&=-2K_{xy;xy},\\
 H_{xy}&=K_{xy;zt}=-K_{xy;tz}, & -2H_{xy}&=+2K_{xy;tz}.
 \label{eq:compare-xy}
\end{align}
两种空间项都写成 \(-2\) 乘各自的原始分量。Unpolarized 的负号
来自 Minkowski 空间指标反序，helicity 的负号来自式
\eqref{eq:helicity-wick-xy}。若直接用 \(F_{tz}\) 写 helicity 的第二端点，
系数为 \(+2\)。\nolinkurl{Helicity_F01_Ft01} 已经用
\(\widetilde F_{01}=F_{23}\) 缩并，因此该文件在式 \eqref{eq:hel} 中
直接乘 \(-2\)。

在同一组态、时间片和距离上先用恒定权重组合各方向，再做位移奇差，
或者先对各张量分量做奇差再组合，结果相同。
式 \eqref{eq:hel} 的差值不除以 2，局域点 \(z=0\) 严格为零。
\(T_H+S_H\) 是空间项反号的对照。

\section{代码、OPE 路径与读取方式}

原始 OPE 根目录
\begin{quote}
\nolinkurl{/public/group/lqcd/donghx/Ope_Gluon/Result_hyp_4D_10times}
\end{quote}
代码位于根目录下的 \nolinkurl{Unified_Ope_code/}。
\nolinkurl{gluon_ope.py} 是 NumPy 入口并包含公共内核。
第 25 行定义空间轮换，第 61 行开始构造 Clover，第 133 行开始取对偶，
第 177 行开始构造 Wilson line 并做缩并。
\nolinkurl{gluon_ope_cupy.py} 是导入该内核的 CuPy/HIP 入口。
输入是已经平滑的规范场，程序直接使用这些链接。
\nolinkurl{Unified_Ope_code/README.md} 记录命令参数、输入的链接轴
\((t,z,y,x,\mathrm{link},c,c)\) 和两个后端入口。

正式文件按
\nolinkurl{<ensemble>/<xdir|ydir|zdir>/<conf>/<filename>.npy}
排列。以下为目录中可读的四个 ensemble。表中的 \(N_t\) 对应输出时间轴。
\begin{center}
\begin{tabular}{lrrr}
\toprule
Ensemble & \(L/a\) & \(N_t\) & \(a\) {\rm [fm]} \\
\midrule
C32P29 & 32 & 64  & 0.1050 \\
E32P29 & 32 & 64  & 0.0897 \\
F32P30 & 32 & 96  & 0.0775 \\
H48P32 & 48 & 144 & 0.0519 \\
\bottomrule
\end{tabular}
\end{center}
文件后缀使用轮换后的实际指标。每行前两个分量相加为 \(T\)，第三个分量乘 2
为 \(S\)，适用于两个通道。
\begin{center}
\begin{tabular}{lccc}
\toprule
Wilson line 方向 & 时间分量 1 & 时间分量 2 & 空间分量 \\
\midrule
\texttt{xdir} & \(31\) & \(32\) & \(12\) \\
\texttt{ydir} & \(32\) & \(30\) & \(20\) \\
\texttt{zdir} & \(30\) & \(31\) & \(01\) \\
\bottomrule
\end{tabular}
\end{center}

以 \nolinkurl{F32P30/zdir/25400/} 为例，\(U_{30}(+z)\) 对应
\nolinkurl{Unpol_F30_F30_dz23.npy}。\(H_{30}(+z)\) 和 \(H_{30}(-z)\) 分别对应
\nolinkurl{Helicity_F30_Ft30_dz23.npy} 和
\nolinkurl{Helicity_F30_Ft30_minus_dz23.npy}。数组为
\texttt{complex128}，轴为 \((t/a,|z|/a)\)，\texttt{dz23} 含
\(0,\ldots,23\) 共 24 列。旧档也有 \texttt{dz19} 和 \texttt{dz29}。
组合时用相同组态、方向、时间片和共有距离列。保留复数到下游投影。
旧的 \nolinkurl{Unpol_2F01F01_minus_F30F30_minus_F31F31}
文件名对应 \(S_U-T_U\)，需要变号才能与式 \eqref{eq:unpol} 对齐。
\nolinkurl{_receipts/} 记录组态级生产或导入来源和文件信息。

\paragraph{落盘 helicity 抽检}
从 C32P29/37800、E32P29/8700、F32P30/1000、H48P32/2480 的
HYP 规范场重建场强，用显式互补指标和正负 Wilson line 复算
\texttt{zdir} 的三个分量。每组检查 \(t/a=0,N_t/2\) 和 \(|z|/a=0,1,2,3\)。
192 个原始分量和 32 个奇差组合的最大绝对残差分别为
\(2.31\times10^{-14}\) 和 \(2.49\times10^{-14}\)。
另取 F32P30/39700 沿相同步骤与旧导入的 helicity 数据比较，
用该组态生产回执中的 HYP 输入复算，原始分量最大绝对差为
\(1.14\times10^{-3}\)，奇差最大绝对差为 \(5.73\times10^{-4}\)。
这是该组态的数值复现差异，具体来源待确认。组态 1000 的通过结果
适用于本次抽检的切片和距离。

\begin{thebibliography}{9}
\bibitem{Khan}
T.~Khan et al., Phys.\ Rev.\ D \textbf{104}, 094516 (2021),
\href{https://arxiv.org/abs/2107.08960v2}{arXiv:2107.08960v2}.
\bibitem{BMRU}
I.~Balitsky, W.~Morris and A.~Radyushkin, Phys.\ Lett.\ B \textbf{808},
135621 (2020), \href{https://arxiv.org/abs/1910.13963}{arXiv:1910.13963}.
\bibitem{Chen}
C.~Chen et al., \emph{Unpolarized gluon PDF of the nucleon from lattice QCD
at physical point in the continuum limit},
\href{https://arxiv.org/abs/2510.26425v2}{arXiv:2510.26425v2}.
\bibitem{Egerer}
C.~Egerer et al., Phys.\ Rev.\ D \textbf{106}, 094511 (2022),
\href{https://arxiv.org/abs/2207.08733v2}{arXiv:2207.08733v2}.
\bibitem{BMRH}
I.~Balitsky, W.~Morris and A.~Radyushkin, JHEP \textbf{02}, 193 (2022),
\href{https://arxiv.org/abs/2112.02011}{arXiv:2112.02011}.
\bibitem{FlowNote}
H.~Dong, \emph{Gradient-Flowed Quark and Gluon Quasi-PDFs in Lattice QCD},
第 8 节，\nolinkurl{/public/home/donghx/Note/Gradient_Flow/main.pdf}.
\end{thebibliography}
\end{document}
